\section{习题}
\label{sec:ch10-exercises}
% ============================================================================

这一章的习题安排得比前几章更像一次总复习。
因为终章真正训练的，
不是某个孤立技巧，
而是把整本书的不同语言在同一个问题上来回切换：
一会儿把指数归约到素数，
一会儿把整数方程送进椭圆曲线，
一会儿再把椭圆曲线送进模形式空间。
下面六题正沿着这条主线展开，
并且每题都配有完整解答。

\begin{figure}[ht]
\centering
\begin{tikzpicture}[>=Stealth, font=\small, node distance=1.05cm]
  \node[draw, rounded corners, fill=blue!8, minimum width=4.8cm, minimum height=0.9cm] (e1) {习题 1：把指数归约到 $4$ 或奇素数};
  \node[draw, rounded corners, fill=green!10, below=of e1, minimum width=4.8cm, minimum height=0.9cm] (e2) {习题 2：用无穷递降解决 $n=4$};
  \node[draw, rounded corners, fill=yellow!16, below=of e2, minimum width=4.8cm, minimum height=0.9cm] (e3) {习题 3：计算 Frey 曲线判别式与最小判别式};
  \node[draw, rounded corners, fill=orange!14, below=of e3, minimum width=4.8cm, minimum height=0.9cm] (e4) {习题 4：验证 $\Sk_2(\GammaO(2))=0$};
  \node[draw, rounded corners, fill=purple!10, below=of e4, minimum width=4.8cm, minimum height=0.9cm] (e5) {习题 5：说明 $X_0(11)\to E_{11}$ 的存在};
  \node[draw, rounded corners, fill=red!8, below=of e5, minimum width=4.8cm, minimum height=0.9cm] (e6) {习题 6：解释为什么是 $3$--$5$ 技巧};

  \draw[->, very thick] (e1) -- (e2);
  \draw[->, very thick] (e2) -- (e3);
  \draw[->, very thick] (e3) -- (e4);
  \draw[->, very thick] (e4) -- (e5);
  \draw[->, very thick] (e5) -- (e6);
\end{tikzpicture}
\caption{终章习题的路线图：从指数归约与无穷递降出发，经 Frey 曲线与空空间 $\Sk_2(\GammaO(2))$，最后回到 Wiles 证明中最精巧的 $3$--$5$ 技巧。}
\label{fig:ch10-exercises-roadmap}
\end{figure}


\begin{exercise}
证明：若费马大定理对 $n=4$ 与所有奇素数 $p$ 成立，
则它对所有 $n\ge 3$ 成立。
\end{exercise}

\begin{proof}[解答]
这正是 \cref{lem:ch10-reduction} 的内容，
但我们在此重写一遍。
设某个 $n\ge 3$ 存在正整数解
\[
  x^n+y^n=z^n.
\]
若 $n$ 有奇素因子 $p$，令
\[
  X=x^{n/p},\qquad Y=y^{n/p},\qquad Z=z^{n/p},
\]
则
\[
  X^p+Y^p=Z^p,
\]
与奇素数情形的 FLT 矛盾。
若 $n$ 没有奇素因子，
则 $n=2^r$ 且 $r\ge 2$。
令
\[
  X=x^{2^{r-2}},\qquad Y=y^{2^{r-2}},\qquad Z=z^{2^{r-2}},
\]
则
\[
  X^4+Y^4=Z^4,
\]
于是得到 $n=4$ 的反例，矛盾。
故只要排除 $n=4$ 与所有奇素数指数，
就排除了全部 $n\ge 3$。
\end{proof}

\begin{exercise}
用 Fermat 的无穷递降法证明
\[
  x^4+y^4=z^4
\]
无正整数解。
更强地，证明
\[
  x^4+y^4=z^2
\]
无正整数解。
\end{exercise}

\begin{proof}[解答]
证明与 \cref{thm:ch10-n4} 完全相同。
反设存在正整数解，并取 $z$ 最小的一组。
约去公因子并交换 $x,y$ 后，可设 $(x,y)=1$ 且 $y$ 偶、$x$ 奇。
于是
\[
  x^2,\ y^2,\ z
\]
构成 primitive Pythagorean triple，故存在互素整数 $m>n>0$，$m$ 奇、$n$ 偶，使得
\[
  x^2=m^2-n^2,
  \qquad y^2=2mn,
  \qquad z=m^2+n^2.
\]
由 $2mn$ 是平方且 $(m,n)=1$，知
\[
  m=r^2,
  \qquad n=2s^2.
\]
于是
\[
  x^2=(r^2-2s^2)(r^2+2s^2).
\]
两因子互素且皆为奇数，乘积是平方，故各自都是平方：
\[
  r^2-2s^2=u^2,
  \qquad r^2+2s^2=v^2.
\]
于是
\[
  (v-u)(v+u)=4s^2.
\]
由 $u,v$ 互素奇数，存在互素整数 $a>b>0$ 使得
\[
  v+u=2a^2,
  \qquad v-u=2b^2.
\]
从而
\[
  u=a^2-b^2,
  \qquad s=ab.
\]
代回可得
\[
  r^2=u^2+2s^2=(a^2-b^2)^2+2a^2b^2=a^4+b^4.
\]
于是 $(a,b,r)$ 给出了更小的正整数解
\[
  a^4+b^4=r^2,
  \qquad r<z,
\]
与 $z$ 最小矛盾。
故 $x^4+y^4=z^2$ 无正整数解。
若 $x^4+y^4=z^4$ 有解，则令 $Z=z^2$ 即得前一种方程的解，矛盾。
\end{proof}

\begin{exercise}
设
\[
  E_{a,b,c}\colon y^2=x(x-a^p)(x+b^p).
\]
验证该整模型的判别式为
\[
  \Delta=16(abc)^{2p},
\]
并说明在标准最小化后可写成
\[
  \Delta_{\min}=2^{-8}(abc)^{2p}.
\]
为什么“奇素数部分是一个完全 $p$ 次幂”这件事在 FLT 路线中如此不自然而又关键？
\end{exercise}

\begin{proof}[解答]
把三根写成
\[
  e_1=0,
  \qquad e_2=a^p,
  \qquad e_3=-b^p.
\]
则 Legendre 型判别式公式给出
\[
  \Delta=16(e_1-e_2)^2(e_1-e_3)^2(e_2-e_3)^2.
\]
代入得到
\[
  \Delta=16(a^p)^2(b^p)^2(a^p+b^p)^2=16(abc)^{2p}.
\]
若 $(a,b,c)=1$ 且 $2\mid b$，
则奇素数处这个模型已经最小，
故对每个奇素数 $q$，
\[
  v_q(\Delta_{\min})=2p\,v_q(abc).
\]
在 $2$ 处经 Tate 算法标准最小化后，
判别式少去一个 $2^{12}$ 因子，
于是
\[
  \Delta_{\min}=2^{-8}(abc)^{2p}.
\]

“不自然”在于：
对一条一般椭圆曲线，
坏素数在判别式里的指数与导子里的指数并不会这样整齐地被同一个大素数 $p$ 整除。
而 Frey 曲线恰恰满足这一极端形状：
导子只记住平方自由部分，
判别式却在每个奇坏素数处带着 $2p$ 倍数指数。
这使得 mod $p$ 表示在这些地方的惯性作用消失，
从而可以应用水平下降。
若没有这个“判别式像 $p$ 次幂、导子却很瘦”的反差，
FLT 证明链就根本无法启动。
\end{proof}

\begin{exercise}
用维数公式验证
\[
  \dim \Sk_2(\GammaO(2))=0.
\]
并说明这为什么是 FLT 证明中的关键一步。
\end{exercise}

\begin{proof}[解答]
对权 $2$，有
\[
  \dim \Sk_2(\GammaO(N))=g(X_0(N)).
\]
因此只需求 $X_0(2)$ 的亏格。
第5章的亏格公式给出
\[
  g(X_0(N))=1+\frac{\mu(N)}{12}-\frac{e_2(N)}{4}-\frac{e_3(N)}{3}-\frac{c_\infty(N)}{2}.
\]
当 $N=2$ 时，
\[
  \mu(2)=2\Bigl(1+\frac12\Bigr)=3,
  \qquad e_2(2)=1,
  \qquad e_3(2)=0,
  \qquad c_\infty(2)=2.
\]
于是
\[
  g(X_0(2))
  =1+\frac{3}{12}-\frac14-0-\frac22
  =1+\frac14-\frac14-1=0.
\]
故
\[
  \dim \Sk_2(\GammaO(2))=0.
\]

这一步在 FLT 证明中之所以关键，
是因为 Ribet 水平下降最后把 Frey 曲线的 residual 表示逼到级别 $2$。
若 $\Sk_2(\GammaO(2))$ 非零，
那么级别 $2$ 仍可能承载那个 residual 特征值系统，
证明就不会终止。
正因为这个空间是零，
“从假设费马解存在出发得到一个级别 $2$ 新形式”才会立刻变成矛盾。
\end{proof}

\begin{exercise}
对 $N=11$，利用
\[
  \dim \Sk_2(\GammaO(11))=1
\]
与第6章的结论，说明 $J_0(11)$ 是一条椭圆曲线，
并写出参数化映射
\[
  X_0(11)\to J_0(11)\cong E_{11}
\]
的存在性论证。
\end{exercise}

\begin{proof}[解答]
因为
\[
  \dim \Sk_2(\GammaO(11))=g(X_0(11))=1,
\]
所以
\[
  \dim J_0(11)=1.
\]
而一维 Abel 簇正是一条椭圆曲线，
故 $J_0(11)$ 是椭圆曲线。
第7章已给出其经典模型
\[
  E_{11}\colon y^2+y=x^3-x^2-10x-20.
\]

接着利用有理尖点 $[\infty]\in X_0(11)(\Q)$。
Abel--Jacobi 映射定义为
\[
  \phi\colon X_0(11)\longrightarrow J_0(11),
  \qquad P\longmapsto [P]-[\infty].
\]
这是定义在 $\Q$ 上的态射，且并非零映射。
由于 $X_0(11)$ 与 $J_0(11)$ 都是亏格 $1$ 曲线，
任何非零态射都必为有限映射；
又因为它把有理点 $[\infty]$ 送到零元，
所以实际上给出一个同源。
在级别 $11$ 这一最简单情形里，
该映射甚至是同构，
于是得到模参数化
\[
  X_0(11)\xrightarrow{\ \sim\ }J_0(11)\cong E_{11}.
\]
这说明“模曲线参数化椭圆曲线”在级别 $11$ 已经是完全显式的事实。
\end{proof}

\begin{exercise}
解释为什么在 Wiles 的证明中，
当 $\rhobar_{E,3}$ 可约时要换用 $\ell=5$，
而不是随便换用 $7$ 或别的素数。
\end{exercise}

\begin{proof}[解答]
关键有两点。

第一，残余模性的起点来自 Langlands--Tunnell 定理，
而这个定理正好适用于二维 mod $3$ 表示：
其射影像落在 $\PGL_2(\F_3)\cong S_4$，
属于可解范围，
从而能先得到残余模性。
如果一开始改用 $\ell=7$，
并没有相应同样强而可直接使用的“mod $7$ 残余模性起点”。
所以证明的第一步天然地偏向 $3$。

第二，当 $\rhobar_{E,3}$ 可约时，
Wiles 需要一个与 $3$ 形成互补的小素数，
并且要能证明对应 residual 表示不可约。
$5$ 正好满足这点：
若半稳定椭圆曲线的 $3$-与 $5$-torsion 都可约，
则曲线会同时带有有理 $3$-与 $5$-isogeny，
从而给出 $15$-isogeny；
这会与 Mazur 的有理同源分类冲突。
因此
\[
  \rhobar_{E,3}\ \text{可约} \Longrightarrow \rhobar_{E,5}\ \text{不可约}.
\]
再加上模曲线 $X(5)$ 的几何足够简单，
可以构造辅助曲线 $A$ 使 $A[5]\cong E[5]$ 而 $A[3]$ 不可约，
整个 $3$--$5$ 技巧才能闭合。

换成 $7$ 并没有这样精确而低维的互补机制：
既缺少像 Langlands--Tunnell 这样的残余模性起点，
也缺少同样紧凑的“若 $3$ 可约则 $7$ 必不可约”与相应辅助模曲线构造。
因此 $5$ 不是随手挑的，
而是与 mod $3$ 起点严丝合缝配套的那个素数。
\end{proof}
